\documentclass[aspectratio=169]{beamer}

% ==== Theme & basic look ====
\usetheme{Madrid}
\usecolortheme{seahorse}
\setbeamertemplate{navigation symbols}{}
\setbeamertemplate{footline}[frame number]


% Packages for mathematics  
\usepackage{amsmath}  
\usepackage{amsfonts}  
\usepackage{amssymb}  
\usepackage{CJKutf8}
% Support for \Tilde macro used throughout
\providecommand{\Tilde}[1]{\tilde{#1}}
\providecommand{\proofname}{Proof}
\newenvironment<>{proofs}[1][\proofname]{%
    \par
    \def\insertproofname{#1.}%
    \usebeamertemplate{proof begin}#2}
  {\usebeamertemplate{proof end}}
\makeatother

% Math shorthand and vector symbols
\newcommand{\grad}{\nabla}
\newcommand{\EE}{\mathbb{E}}
\newcommand{\dd}{\mathrm{d}}
\newcommand{\kbT}{k_B T}
\newcommand{\vx}{\mathbf{x}}

% Bold vector/time-indexed variables
\newcommand{\vxz}{\mathbf{x}_0}
\newcommand{\vxt}{\mathbf{x}_t}
\newcommand{\alpt}{\alpha_t}
\newcommand{\sigt}{\sigma_t}
\newcommand{\veps}{\boldsymbol{\epsilon}}
\newcommand{\vF}{\mathbf{F}}

% Agent-specific notation
\newcommand{\agent}{\mathcal{A}}
\newcommand{\state}{\mathbf{s}}

% Title page details:   
\title[LLM Agents]{Agent Technologies in Large Language Models: Intelligent Systems Beyond Traditional AI}  
\author{Pipi Hu}  
\institute[BIMSA]{Beijing Institute of Mathematical Sciences and Applications}  
\date{\today}  

\pgfdeclareimage[width=2cm]{logo}{../figures/logo.png}
\logo{\pgfuseimage{logo}}

% Outline slide (moved to later; avoid frame before \begin{document})

\setbeamertemplate{navigation symbols}{}
\AtBeginSection[]  
{  
    \begin{frame}<beamer>{Outline}         
        \tableofcontents[currentsection]  
    \end{frame}  
}  
  
\AtBeginSubsection[]  
{  
    \begin{frame}<beamer>{Outline}         
        \tableofcontents[currentsection,currentsubsection]  
    \end{frame}  
}  

\begin{document}  
  
% Title slide  
\begin{frame}  
  \titlepage  
\end{frame}  

\begin{frame}{The Evolution of AI Agents}
\centering
    "The question of whether a computer can think is like asking whether a submarine can swim."
    \\[1em]
    \textit{--- Edsger W. Dijkstra}
    \\[1em]
    Today's LLM agents can reason, plan, and interact with the world.
\end{frame}
  
% Presentation slides  

\section{Introduction: LLM Agents vs Traditional AI}

\begin{frame}{What are LLM Agents?}
LLM agents are intelligent systems built on large language models that can:
\begin{itemize}
    \item \textbf{Natural Language Processing}: Understand and generate human-like text
    \item \textbf{Tool Integration}: Use external APIs and functions to perform tasks
    \item \textbf{Reasoning}: Perform logical inference and problem-solving
    \item \textbf{Planning}: Create and execute multi-step plans
    \item \textbf{Context Awareness}: Maintain and utilize conversation context
\end{itemize}
\end{frame}

\begin{frame}{LLM Agents vs Traditional AI}
\begin{center}
\begin{tabular}{|l|l|l|}
\hline
\textbf{Aspect} & \textbf{Traditional AI} & \textbf{LLM Agents} \\
\hline
Learning & Supervised/Unsupervised & Few-shot/Prompt-based \\
\hline
Interaction & Rule-based & Natural language \\
\hline
Flexibility & Domain-specific & General-purpose \\
\hline
Tool Use & Limited & Extensive API integration \\
\hline
Reasoning & Pre-programmed & Emergent from training \\
\hline
\end{tabular}
\end{center}
\end{frame}

\section{Mathematical Foundations of LLM Agents}

\begin{frame}{LLM Agent Architecture}
An LLM agent can be modeled as:
\begin{equation}
    \agent_{LLM}: \mathcal{P} \times \mathcal{C} \times \mathcal{T} \rightarrow \mathcal{R}
\end{equation}
where:
\begin{itemize}
    \item $\mathcal{P}$: Prompt space (user inputs)
    \item $\mathcal{C}$: Context space (conversation history)
    \item $\mathcal{T}$: Tool space (available functions/APIs)
    \item $\mathcal{R}$: Response space (text outputs)
\end{itemize}
\end{frame}

\begin{frame}{LLM Agent Decision Process}
The agent's response generation:
\begin{equation}
    \text{Response} = \text{LLM}(\text{Prompt} \| \text{Context} \| \text{Tool\_Calls})
\end{equation}

The complete process:
\begin{align}
    \text{Understanding} &= \text{Parse}(\text{Prompt}) \\
    \text{Planning} &= \text{Reason}(\text{Understanding}, \text{Context}) \\
    \text{Execution} &= \text{Act}(\text{Planning}, \text{Tools}) \\
    \text{Response} &= \text{Generate}(\text{Execution\_Results})
\end{align}
\end{frame}

\begin{frame}{Multi-LLM Agent Collaboration}
For a system with $N$ LLM agents $\{\agent_1, \agent_2, \ldots, \agent_N\}$:
\begin{equation}
    \text{Output} = \text{Coordination}(\agent_1(\text{Task}_1), \agent_2(\text{Task}_2), \ldots, \agent_N(\text{Task}_N))
\end{equation}

Each agent's response:
\begin{equation}
    \text{Response}_i = \agent_i(\text{Prompt}_i, \text{Context}_i, \text{Tools}_i)
\end{equation}

Inter-agent communication:
\begin{equation}
    \text{Context}_i = \text{Context}_i \cup \{\text{Response}_j : j \neq i\}
\end{equation}
\end{frame}

\begin{frame}{LLM Agent Prompt Engineering}
The effectiveness of an LLM agent depends on prompt quality:
\begin{equation}
    \text{Performance} = f(\text{Prompt\_Quality}, \text{Context\_Relevance}, \text{Tool\_Availability})
\end{equation}

Prompt optimization:
\begin{equation}
    \text{Optimal\_Prompt} = \arg\max_{\text{Prompt}} \text{Task\_Success\_Rate}(\text{Prompt})
\end{equation}

Chain-of-thought reasoning:
\begin{equation}
    \text{Response} = \text{LLM}(\text{Problem} + \text{"Let's think step by step..."})
\end{equation}
\end{frame}

\section{LLM Agent Architectures}

\begin{frame}{Core Components of LLM Agents}
LLM agents integrate several key components:
\begin{enumerate}
    \item \textbf{Language Model}: Core transformer architecture for text understanding
    \item \textbf{Function Calling}: Ability to invoke external tools and APIs
    \item \textbf{Memory System}: Conversation history and context management
    \item \textbf{Reasoning Engine}: Chain-of-thought and logical inference
    \item \textbf{Planning Module}: Multi-step task decomposition and execution
\end{enumerate}
\end{frame}

\begin{frame}{LLM Agent Processing Pipeline}
The LLM agent processes tasks through a structured pipeline:
\begin{equation}
    \text{Output} = \text{LLM}(\text{System\_Prompt} \| \text{User\_Input} \| \text{Context} \| \text{Tool\_Results})
\end{equation}

The processing stages:
\begin{align}
    \text{Parse} &: \text{Input} \rightarrow \text{Structured\_Intent} \\
    \text{Plan} &: \text{Intent} \rightarrow \text{Action\_Sequence} \\
    \text{Execute} &: \text{Sequence} \rightarrow \text{Tool\_Calls} \\
    \text{Synthesize} &: \text{Results} \rightarrow \text{Final\_Response}
\end{align}
\end{frame}

\begin{frame}{Function Calling in LLM Agents}
LLM agents can call external functions through structured outputs:
\begin{equation}
    \text{Function\_Call} = \text{LLM}(\text{Input}) \rightarrow \{\text{function\_name}, \text{parameters}\}
\end{equation}

Function execution:
\begin{equation}
    \text{Result} = \text{Execute}(\text{Function\_Call}) \rightarrow \text{Response}
\end{equation}

The complete cycle:
\begin{equation}
    \text{User\_Input} \rightarrow \text{LLM} \rightarrow \text{Function\_Call} \rightarrow \text{Result} \rightarrow \text{Response}
\end{equation}
\end{frame}

\section{Multi-LLM Agent Systems}

\begin{frame}{LLM Agent Collaboration Patterns}
Multiple LLM agents can work together in different patterns:
\begin{enumerate}
    \item \textbf{Sequential}: $\agent_1 \rightarrow \agent_2 \rightarrow \agent_3$
    \item \textbf{Parallel}: $\agent_1 \| \agent_2 \| \agent_3$
    \item \textbf{Hierarchical}: $\agent_{master} \rightarrow \{\agent_{worker_i}\}$
    \item \textbf{Democratic}: $\text{Vote}(\{\agent_i\})$
\end{enumerate}
\end{frame}

\begin{frame}{Agent Communication Protocols}
LLM agents communicate through structured messages:
\begin{equation}
    \text{Message} = \{\text{sender}, \text{receiver}, \text{content}, \text{timestamp}\}
\end{equation}

Message routing:
\begin{equation}
    \text{Route}(\text{Message}) = \text{Find\_Agent}(\text{receiver})
\end{equation}
\end{frame}

\begin{frame}{LLM Agent Specialization}
Different LLM agents can specialize in different tasks:
\begin{itemize}
    \item \textbf{Research Agent}: Information gathering and analysis
    \item \textbf{Coding Agent}: Software development and debugging
    \item \textbf{Creative Agent}: Content generation and design
    \item \textbf{Analytical Agent}: Data processing and insights
    \item \textbf{Communication Agent}: Human interaction and translation
\end{itemize}
\end{frame}

\section{Tool Integration in LLM Agents}

\begin{frame}{Function Calling Architecture}
LLM agents can integrate with external tools through function calling:
\begin{equation}
    \text{Tool\_Call} = \text{LLM}(\text{Input}) \rightarrow \{\text{name}, \text{arguments}\}
\end{equation}

Tool execution pipeline:
\begin{align}
    \text{Parse} &: \text{LLM\_Output} \rightarrow \text{Function\_Spec} \\
    \text{Validate} &: \text{Spec} \rightarrow \text{Validated\_Call} \\
    \text{Execute} &: \text{Call} \rightarrow \text{Result} \\
    \text{Format} &: \text{Result} \rightarrow \text{Response}
\end{align}
\end{frame}

\begin{frame}{Common Tool Categories for LLM Agents}
LLM agents can integrate with various tool categories:
\begin{enumerate}
    \item \textbf{Web Search}: Real-time information retrieval
    \item \textbf{Code Execution}: Running and debugging code
    \item \textbf{Database Access}: Querying structured data
    \item \textbf{File Operations}: Reading and writing files
    \item \textbf{API Integration}: Third-party service calls
\end{enumerate}
\end{frame}

\section{Reasoning and Planning in LLM Agents}

\begin{frame}{Chain-of-Thought Reasoning}
LLM agents use chain-of-thought prompting for complex reasoning:
\begin{equation}
    \text{Response} = \text{LLM}(\text{Problem} + \text{"Let's think step by step..."})
\end{equation}

The reasoning process:
\begin{align}
    \text{Step}_1 &: \text{Understand}(\text{problem}) \rightarrow \text{key\_concepts} \\
    \text{Step}_2 &: \text{Decompose}(\text{problem}) \rightarrow \text{subproblems} \\
    \text{Step}_3 &: \text{Solve}(\text{subproblems}) \rightarrow \text{solutions} \\
    \text{Step}_4 &: \text{Synthesize}(\text{solutions}) \rightarrow \text{final\_answer}
\end{align}
\end{frame}

\begin{frame}{Task Planning in LLM Agents}
LLM agents break down complex tasks into manageable steps:
\begin{equation}
    \text{Task\_Plan} = \text{LLM}(\text{Goal} + \text{"Break this down into steps"})
\end{equation}

Planning components:
\begin{align}
    \text{Goal\_Analysis} &: \text{Understand}(\text{objective}) \\
    \text{Step\_Generation} &: \text{Decompose}(\text{goal}) \\
    \text{Sequence\_Ordering} &: \text{Arrange}(\text{steps}) \\
    \text{Execution\_Plan} &: \text{Optimize}(\text{sequence})
\end{align}
\end{frame}

\section{Memory and Context in LLM Agents}

\begin{frame}{Context Window Management}
LLM agents manage conversation context within token limits:
\begin{equation}
    \text{Context} = \text{System\_Prompt} + \text{Conversation\_History} + \text{Current\_Input}
\end{equation}

Context optimization:
\begin{equation}
    \text{Optimized\_Context} = \text{Truncate}(\text{Context}, \text{max\_tokens})
\end{equation}
\end{frame}

\begin{frame}{Long-Term Memory for LLM Agents}
LLM agents can maintain long-term memory through:
\begin{itemize}
    \item \textbf{Vector Databases}: Semantic similarity search
    \item \textbf{External Storage}: Persistent conversation history
    \item \textbf{Knowledge Graphs}: Structured relationship storage
    \item \textbf{Embedding Caches}: Pre-computed similarity vectors
\end{itemize}
\end{frame}

\section{LLM Agent Applications}

\begin{frame}{Real-World LLM Agent Applications}
LLM agents are being deployed in various domains:
\begin{itemize}
    \item \textbf{Code Generation}: GitHub Copilot, ChatGPT Code Interpreter
    \item \textbf{Research Assistance}: Literature review and data analysis
    \item \textbf{Customer Service}: Automated support and chatbots
    \item \textbf{Content Creation}: Writing, editing, and creative tasks
    \item \textbf{Data Analysis}: Business intelligence and insights
\end{itemize}
\end{frame}

\begin{frame}{LLM Agent Success Stories}
Notable examples of successful LLM agent implementations:
\begin{enumerate}
    \item \textbf{AutoGPT}: Autonomous task execution and planning
    \item \textbf{LangChain Agents}: Modular agent framework
    \item \textbf{ReAct Framework}: Reasoning and acting in language models
    \item \textbf{Toolformer}: Teaching LLMs to use tools
    \item \textbf{WebGPT}: Web-based question answering
\end{enumerate}
\end{frame}

\section{Evaluation of LLM Agents}

\begin{frame}{LLM Agent Evaluation Metrics}
Key metrics for evaluating LLM agent performance:
\begin{align}
    \text{Task\_Success\_Rate} &= \frac{\text{Completed\_Tasks}}{\text{Total\_Tasks}} \\
    \text{Response\_Quality} &= \text{Human\_Rating}(\text{outputs}) \\
    \text{Tool\_Usage\_Accuracy} &= \frac{\text{Correct\_Tool\_Calls}}{\text{Total\_Tool\_Calls}} \\
    \text{Context\_Retention} &= \frac{\text{Relevant\_Responses}}{\text{Total\_Responses}}
\end{align}
\end{frame}

\begin{frame}{Benchmarking LLM Agents}
Standard benchmarks for evaluating LLM agents:
\begin{itemize}
    \item \textbf{WebArena}: Web-based task automation
    \item \textbf{AgentBench}: Multi-domain agent evaluation
    \item \textbf{ALFWorld}: Interactive household tasks
    \item \textbf{ScienceQA}: Scientific reasoning tasks
    \item \textbf{HotpotQA}: Multi-hop reasoning
\end{itemize}
\end{frame}

\section{Challenges and Future Directions}

\begin{frame}{Current Challenges in LLM Agents}
\begin{enumerate}
    \item \textbf{Hallucination}: Generating factually incorrect information
    \item \textbf{Context Limits}: Managing long conversations within token limits
    \item \textbf{Tool Reliability}: Ensuring consistent external API usage
    \item \textbf{Safety}: Preventing harmful or biased outputs
    \item \textbf{Cost}: Managing computational expenses for large-scale deployment
\end{enumerate}
\end{frame}

\begin{frame}{Future Research Directions}
\begin{itemize}
    \item \textbf{Multimodal Agents}: Integration of text, vision, and audio
    \item \textbf{Longer Context Windows}: Handling extended conversations
    \item \textbf{Specialized Agents}: Domain-specific expertise
    \item \textbf{Agent Swarms}: Large-scale collaborative systems
    \item \textbf{Constitutional AI}: Built-in ethical guidelines
\end{itemize}
\end{frame}

\section{Implementation Considerations}

\begin{frame}{LLM Agent System Architecture}
Key components for deploying LLM agents:
\begin{enumerate}
    \item \textbf{LLM Backend}: Core language model (GPT, Claude, etc.)
    \item \textbf{Function Registry}: Available tools and APIs
    \item \textbf{Context Manager}: Conversation history and memory
    \item \textbf{Orchestrator}: Task planning and execution
    \item \textbf{Monitoring}: Performance tracking and logging
\end{enumerate}
\end{frame}

\begin{frame}{LLM Agent Development Best Practices}
\begin{itemize}
    \item \textbf{Prompt Engineering}: Clear, specific instructions
    \item \textbf{Function Design}: Well-defined tool interfaces
    \item \textbf{Error Handling}: Graceful failure recovery
    \item \textbf{Testing}: Comprehensive evaluation on diverse tasks
    \item \textbf{Safety}: Built-in safeguards and monitoring
\end{itemize}
\end{frame}

\section{Conclusion}

\begin{frame}{Key Takeaways}
LLM agents represent a paradigm shift in AI capabilities:
\begin{enumerate}
    \item \textbf{Natural Language Interface}: Human-like communication
    \item \textbf{Tool Integration}: Seamless external system interaction
    \item \textbf{Reasoning}: Chain-of-thought problem solving
    \item \textbf{Planning}: Multi-step task execution
    \item \textbf{Adaptability}: Few-shot learning and context awareness
\end{enumerate}
\end{frame}

\begin{frame}{The Future of LLM Agents}
As we look forward:
\begin{itemize}
    \item \textbf{Multimodal Capabilities}: Text, image, and audio integration
    \item \textbf{Longer Context}: Extended conversation memory
    \item \textbf{Specialized Agents}: Domain-specific expertise
    \item \textbf{Agent Ecosystems}: Large-scale collaborative networks
    \item \textbf{Ethical AI}: Built-in safety and alignment
\end{itemize}
\end{frame}

\begin{frame}{Questions and Discussion}
\begin{center}
\textbf{Key Questions for Further Research:}
\begin{enumerate}
    \item How can we reduce hallucination in LLM agents?
    \item What are the optimal prompt engineering strategies?
    \item How can we improve tool usage reliability?
    \item What metrics best evaluate LLM agent performance?
    \item How can we ensure safety in autonomous LLM agents?
\end{enumerate}
\end{center}
\end{frame}

% Thank you slide  
\begin{frame}{Welcome inputs}  
  \centering  
  Any comments?
  \\[2em]
  Welcome your inputs!  
\end{frame}  
  
\end{document}  
